% =====================================================================
\section{Exercise 5 -- Berger (2016) simulation: Kupiec backtesting under
copula misspecification}\label{sec:ex5}
% =====================================================================

We replicate Section~3.1 of Berger (2016) restricted to a Clayton DGP. The
goal is to quantify how copula misspecification distorts VaR coverage. The
implementation is in \file{exercise5/} (driver \file{exercise5/main.py}).

% ---------------------------------------------------------------
\subsection{Simulation design}
% ---------------------------------------------------------------

Each of $R$ Monte-Carlo replications proceeds in five steps:

\begin{enumerate}\setlength\itemsep{0pt}
\item \textbf{DGP.} Draw $T+1=1001$ pairs $(U_{1t},U_{2t})$ from a Clayton
      copula with $\theta\in\{0.5,1.5,2.5\}$ and set
      $r_{it}=\Phi^{-1}(U_{it})$.
\item \textbf{In-sample fit.} On observations $1\ldots T=1000$ estimate
      three competing copula families (Gaussian, $t$, Clayton) treating the
      $\mathcal{N}(0,1)$ marginals as known.
\item \textbf{VaR forecast.} Simulate $M=10{,}000$ pairs from each fitted
      copula, form $r_p^{(s)}=\tfrac12 r_1^{(s)}+\tfrac12 r_2^{(s)}$ and
      record the empirical $\alpha$-quantile
      $\widehat{\mathrm{VaR}}_\alpha$, $\alpha\in\{0.01,0.05\}$ (Berger eq.~9).
\item \textbf{Backtest.} Compare $r_{p,T+1}=\tfrac12 r_{1,T+1}+\tfrac12 r_{2,T+1}$
      with $\widehat{\mathrm{VaR}}_\alpha$ to obtain a single Bernoulli breach
      indicator per replication.
\item \textbf{Aggregate.} After $R=1000$ replications, the empirical failure
      rate is $\hat\pi=N/R$ where $N=\sum_r I_r$, and the Kupiec
      unconditional-coverage statistic (Berger eq.~11) is
      \begin{equation}\label{eq:kupiec}
      LR_{\text{UC}}=-2\ln\bigl[(1-\alpha)^{R-N}\alpha^N\bigr]
                    +2\ln\bigl[(1-\hat\pi)^{R-N}\hat\pi^N\bigr]\;\stackrel{H_0}{\sim}\;\chi^2_1.
      \end{equation}
\end{enumerate}

\paragraph{Estimation details.}
Marginals are fixed at $\mathcal{N}(0,1)$, so the IFM step collapses to a
direct copula MLE on $u_{it}=\Phi(r_{it})$. The Gaussian copula uses
\eqref{eq:gauss-cop-ll} (Brent), the $t$-copula uses \eqref{eq:t-cop-ll}
with the reparameterisation $\rho=\tanh\xi$, $\nu=2+e^\lambda$ (unconstrained
Nelder--Mead). For the Clayton copula we use the closed-form bivariate density
\[
\ln c^{Cla}(u_1,u_2;\theta)=\ln(1+\theta)-(1+\theta)(\ln u_1+\ln u_2)
-\Bigl(2+\tfrac{1}{\theta}\Bigr)\ln\bigl(u_1^{-\theta}+u_2^{-\theta}-1\bigr),
\]
maximised by bounded Brent on $(0.01,30)$.

% ---------------------------------------------------------------
\subsection{Average estimated parameters}
% ---------------------------------------------------------------

Table~\ref{tab:5-params} reports, for each $\theta_{\text{true}}$, the
average of the in-sample fitted parameters across the $R=1000$ replications.

\begin{table}[H]\centering\small
\caption{Average fitted copula parameters across $R=1000$ replications.}
\label{tab:5-params}
\begin{tabular}{c|cc|ccc|cc}
\toprule
& \multicolumn{2}{c|}{Gaussian} & \multicolumn{3}{c|}{Student $t$} & \multicolumn{2}{c}{Clayton}\\
$\theta_{\text{true}}$ & $\hat\rho$ (sd) & $\ln L$ & $\hat\rho$ & $\hat\nu$ & $\ln L$ & $\hat\theta$ (sd) & $\ln L$\\\midrule
$0.5$ & $0.317\ (0.028)$ & --- & $0.315$ & $\to\infty$ & --- & $0.501\ (0.047)$ & ---\\
$1.5$ & $0.609\ (0.020)$ & --- & $0.616$ & $7.34$      & --- & $1.499\ (0.076)$ & ---\\
$2.5$ & $0.736\ (0.016)$ & --- & $0.750$ & $5.10$      & --- & $2.501\ (0.101)$ & ---\\
\bottomrule
\end{tabular}
\end{table}

The Clayton estimator is unbiased and increasingly precise (in absolute terms)
as the dependence strengthens. The $t$-copula degenerates to Gaussian
($\hat\nu\to\infty$) at low Clayton dependence -- the lower-tail signal is
too weak to identify $\nu$ -- but settles at $\hat\nu\approx 5$ at $\theta=2.5$,
consistent with the Clayton lower-tail coefficient
$\lambda_L=2^{-1/\theta}=0.757$ at $\theta=2.5$.

% ---------------------------------------------------------------
\subsection{Empirical VaR coverage and Kupiec test}
% ---------------------------------------------------------------

Tables~\ref{tab:5-95} and \ref{tab:5-99} report the empirical breach rates
together with Kupiec rejection markers ($^*$ at the 95\% level, $^{**}$ at
99\%). Figure~\ref{fig:5} visualises the same numbers as a grouped bar chart.

\begin{table}[H]\centering\small
\caption{Empirical breach rates (\%) at the 95\% VaR ($\alpha=5\%$).
$^*$ Kupiec reject at 5\%; $^{**}$ at 1\%.}\label{tab:5-95}
\begin{tabular}{ccccc}
\toprule
DGP $\theta$ & Gaussian & $t$ & Clayton & Nominal\\\midrule
$0.5$ & $4.60$ & $4.90$ & $4.20$ & $5.00$\\
$1.5$ & $7.40^{**}$ & $7.20^{**}$ & $5.60$ & $5.00$\\
$2.5$ & $6.80^{*}$  & $6.80^{*}$  & $5.70$ & $5.00$\\
\bottomrule
\end{tabular}
\end{table}

\begin{table}[H]\centering\small
\caption{Empirical breach rates (\%) at the 99\% VaR ($\alpha=1\%$).
$^*$ Kupiec reject at 5\%; $^{**}$ at 1\%.}\label{tab:5-99}
\begin{tabular}{ccccc}
\toprule
DGP $\theta$ & Gaussian & $t$ & Clayton & Nominal\\\midrule
$0.5$ & $1.70^{*}$ & $1.80^{*}$ & $1.10$       & $1.00$\\
$1.5$ & $1.70^{*}$ & $1.30$     & $0.90$       & $1.00$\\
$2.5$ & $1.90^{*}$ & $1.90^{*}$ & $1.70^{*}$   & $1.00$\\
\bottomrule
\end{tabular}
\end{table}

\begin{figure}[H]\centering
\includegraphics[width=0.95\linewidth]{fig_5_failure_rates.png}
\caption{Empirical failure rates by DGP intensity, copula family, and confidence
level. Dashed line is the nominal $\alpha$.}
\label{fig:5}
\end{figure}

% ---------------------------------------------------------------
\subsection{Reading the table -- Berger's central message}
% ---------------------------------------------------------------

\begin{enumerate}\setlength\itemsep{0pt}
\item At low Clayton dependence ($\theta=0.5$, Kendall's $\tau=0.20$) the
      three copulas perform similarly at 95\%, but Gaussian and $t$ already
      reject Kupiec at the 99\% level: \emph{the lower tail is too thin under
      a Gaussian copula}, breaches occur 1.7--1.8\% of the time vs.\ a
      nominal 1\%.
\item At medium dependence ($\theta=1.5$, $\tau=0.43$) the misspecification
      is severe at the 95\% level: Gaussian 7.4\% (\textbf{**}), $t$ 7.2\%
      (\textbf{**}). The $t$-copula partially repairs the 99\% miscoverage
      because it matches the heavier joint tails, but Gaussian remains
      rejected.
\item At high dependence ($\theta=2.5$, $\tau=0.56$) all three models start
      under-forecasting losses at 99\%; even the correctly-specified Clayton
      drifts up to 1.7\% (rejected at 5\%) due to small-sample VaR-quantile
      noise.
\end{enumerate}

The conclusion aligned with Berger: \textbf{ignoring lower-tail
dependence biases VaR forecasts toward optimism}. A dynamic risk-management
system anchored on a Gaussian copula will systematically underestimate
losses, and the unconditional-coverage Kupiec test reliably detects this
under the simulation lengths used in practice.

% ---------------------------------------------------------------
\subsection{Implementation notes}
% ---------------------------------------------------------------

Compared to the original paper, our \file{config.json} uses $R=1000$
replications instead of $R=10{,}000$ to reduce computation time; the
qualitative ranking of the three copulas is preserved, but this increases
the standard error.
